% total-accumulation.tex
\documentclass[border=10pt]{standalone}
\usepackage{pgfplots}
\usepackage{pgfplotstable}
\pgfplotsset{compat=1.18}

% ---- USER SETTINGS ----
\def\startday{0}      % change to e.g. 5 to start from day 5 (baseline = volume at day 4)
\def\stopday{97}      % change to e.g. 60 to stop at day 60 (large number = all data)

\begin{document}
% Read CSV with explicit comma separator
\pgfplotstableread[col sep=comma]{experimentos-day.csv}{\data}

% Determine baseline volume (value at day = startday - 1)
\ifnum\startday>0
    \pgfplotstablegetelem{\the\numexpr\startday-1\relax}{[index]1}\of\data
    \edef\baseline{\pgfplotsretval}
\else
    \def\baseline{0}
\fi

\begin{tikzpicture}
\begin{axis}[
    width=15cm, height=10cm,
    ybar,
    bar width=0.8,
    xlabel={Day},
    ylabel={Volume (NmL)},
    title={Cumulative Volume\ifnum\startday>0\ (relative to day \the\numexpr\startday-1\relax)\fi},
    xmin=\startday, xmax=\stopday,
    ymin=0,
    grid=major,
    xtick distance=10,
    minor xtick={1,2,...,\stopday},
    x tick label style={/pgf/number format/1000 sep=},
    restrict x to domain=\startday:\stopday,
]
\addplot[
    fill=blue!50,
    draw=blue,
] table[x index=0, y expr={\thisrowno{1} - \baseline}] {\data};
\end{axis}
\end{tikzpicture}
\end{document}
